\subsection{An all-orders quadratic Gauss obstruction}
\label{subsec:quadratic-gauss-family}
The trace-ideal reduction also gives a complete exclusion at
multiplicity eight.  The argument below treats every possible central
restriction column and every root-of-unity order; it does not impose a
numerical conductor cutoff.

Write $R(m)$ for the commutative self-dual rank-six tensor on
$1,a,b,c,d,e$ whose nonunit coefficients equal $m$, except that the four
unordered triples
\[
 (a,a,b),\ (a,a,c),\ (a,b,b),\ (b,b,c)
\]
have coefficient $m-1$, and
\[
 (c,d,d),\ (c,e,e),\ (d,d,e),\ (d,e,e)
\]
have coefficient $m+1$.  Coefficients are invariant under permutations of
the three indices, and the unit coefficients are fixed by self-duality.
Set
\[
 n=25m^2+4,\quad U=\frac{5m+\sqrt n}{2},\quad
 j=(1,1,1,1,1)^t,\quad\delta=(-1,-1,0,1,1)^t.
\]
The FP dimensions are $(1,U-1,U-1,U,U+1,U+1)$, and
$D=5(U^2+1)=(5n+25m\sqrt n)/2$.
The other five characters are the quadratic conjugate of the FP
character and, with $p=(1+\sqrt5)/2$, $q=(1-\sqrt5)/2$,
\[
 (1,p,q,-1,0,0),\ (1,q,p,-1,0,0),\
 (1,0,0,1,-p,-q),\ (1,0,0,1,-q,-p).
\]
The corresponding formal codegrees are $D,D',5,5,5,5$.
These identities, character completeness, and the full tensor are checked
exactly by \filename{software/new\_gauss\_family.py}.

\begin{proposition}[Quadratic Gauss family obstruction]
\label{prop:quadratic-gauss-family}
The tensor $R(7)=\RR{8}{6520}$ has no pseudo-unitary categorification.
Consequently the trace-ideal signed reduction excludes every complex
fusion categorification of this tensor.
More generally, the same exact verifier gives a necessary criterion for
$R(m)$, $m\ge3$: whenever its exhaustive column enumeration has the
families described below, it requires
\[
 12m^2+2\ \ge\ \frac{\varphi(f)m b}{2c},\qquad
 n=b^2q_0,\quad f=\operatorname{disc}\QQ(\sqrt{q_0}),\quad
 c=\sqrt{f/q_0},
\]
where $q_0$ is squarefree.
\end{proposition}
\begin{proof}
Assume a pseudo-unitary categorification and use its positive spherical
structure.  If $v_Y$ is the restriction column of a central simple, then
$\sum_Yv_Yv_Y^t=G$, where $G=\sum_iN_i^2$.
The unit column of $G$ is $(6,5mj+\delta)^t$, and its nonunit block is
\[
 5I+(25m^2+1)jj^t+5m(j\delta^t+\delta j^t)+\delta\delta^t.
\]
By central induction, $I(1)$ contains six distinct simples with dimensions
$1,z=U^2,t=U^2+1,t,t,t$ and twists one
\cite[Theorems 2.5 and 2.13]{Ostrik3}.
The symmetry of the central $S$-matrix gives, for character-labelled low
columns,
\[
 \dim(Y_\chi)\chi(v_\psi)=\dim(Y_\psi)\psi(v_\chi).
\]
Applying this with the quadratic-conjugate character and the four
codegree-five characters determines the $z$ column as $(1,mj)^t$:
indeed the conjugate dimension of a $t$ column is $t'=t/z$, so
reciprocity gives each small character value as $zt'/t=1$.  The complete
character table then determines the column uniquely.  Write the remaining four columns as
$(1,mj+w_\nu)^t$.  Their dimensions and the unit column of the Gram identity
give
\[
 j^tw_\nu=0,\qquad \delta^tw_\nu=1,\qquad
 \sum_\nu w_\nu=\delta.
\]
On the three-dimensional subspace perpendicular to $j$ and $\delta$,
$G$ restricts to $5I$.  Taking the trace of its residual Gram inequality
therefore gives $\sum_\nu\|w_\nu\|^2\le16$.  Each squared norm is at least
two, so every $w_\nu$ has squared norm at most ten and coordinates in
$[-3,3]$.

The exhaustive enumeration has 32 individual deviations, 59 unordered
quadruples satisfying the sum and norm conditions, and 11 satisfying a
labelling by the four small characters and modular reciprocity.
Exactly three have positive-semidefinite residual Gram matrix.  Their
rows of deviations are
\[
\begin{split}
W_1&=\begin{pmatrix}-1&0&1&-1&1\\-1&1&-1&1&0\\0&-1&1&1&-1\\1&-1&-1&0&1\end{pmatrix},\qquad
W_2=\begin{pmatrix}-1&0&1&0&0\\0&-1&1&0&0\\0&0&-1&0&1\\0&0&-1&1&0\end{pmatrix},\\
W_3&=\begin{pmatrix}-1&0&1&1&-1\\-1&1&-1&0&1\\0&-1&1&-1&1\\1&-1&-1&1&0\end{pmatrix}.
\end{split}
\]
All enumeration bounds have just been proved; no low column is selected
in advance.

A remaining column is $(0,v)^t$, with $v\in\ZZ_{\ge0}^5\setminus\{0\}$.
Let $Q$ be the nonunit residual Gram matrix.  The necessary conditions are
\[
 Q-vv^t\succeq0,\qquad Q-vv^t\ge0\text{ entrywise},\qquad
 D/\dim(v)\text{ algebraically integral}.
\]
The last condition follows because $D^2/\dim(v)^2$ is a formal codegree
of the modular center.  All three residuals annihilate
$(1,1,-4,1,1)^t$.  Hence the first, second, fourth, and fifth coordinates
determine the third.  The bounds $v_i\le\lfloor\sqrt{Q_{ii}}\rfloor$ and
\[
 |v_a-v_b|\le\lfloor\sqrt{Q_{aa}+Q_{bb}-2Q_{ab}}\rfloor,
 \qquad
 |v_d-v_e|\le\lfloor\sqrt{Q_{dd}+Q_{ee}-2Q_{de}}\rfloor
\]
make the enumeration finite and complete.  Gram positivity is checked
with exact rational Schur complements, or equivalently by the kernel
condition and $v_J^tQ_J^{-1}v_J\le1$ on a full-rank principal submatrix.
For a quadratic number $x+y\sqrt n$, integrality is equivalent to integral
trace $2x$ and integral norm $x^2-ny^2$.

For $m=7$ the counts at the bounded, Gram-compatible, and integral-quotient
stages are respectively
\[
 (511,58,1),\qquad (6147,282,9),\qquad (511,58,1).
\]
The outer branches admit only $j$, although $Q$ has rank two, and are
therefore impossible.  The complete middle list is
\[
 j,\quad j\pm(e_a-e_b),\quad j\pm(e_d-e_e),\quad
 7j-e_i+e_k\quad(i\in\{a,b\},\ k\in\{d,e\}).
\]
For the general verifier, the accepted list may also include the first
five types with $j$ replaced by $2j$ and the last four types with $m$
replaced by $2m$; any additional high type causes the criterion to abstain.

Put $u=e_a-e_b$, $v=e_d-e_e$.  Exact Gram identities are
\[
 \delta^tQ\delta=32,\quad\delta^tQj=80m,\quad
 u^tQu=v^tQv=8,\quad
 \delta^tQu=\delta^tQv=u^tQv=0.
\]
If $n_m,n_{2m}$ count the two possible cross-column families, the first
two equations give $n_m+n_{2m}=8$ and $n_m+2n_{2m}=8$.
Thus exactly eight cross columns occur, all with coefficient $m$.
They saturate both $u,v$ squared forms, excluding all offset types.
The three vanishing mixed forms force each of the four columns
$B_{ik}=mj-e_i+e_k$ to occur twice.  Subtracting them leaves
\[
 Q-2\sum_{i,k}B_{ik}B_{ik}^t=(12m^2+2)jj^t.
\]
Thus the number of uniform columns, with a column $2j$ counted four
times, is $L=12m^2+2$.

Let $\Theta_{ik}$ be the sum of the two twists belonging to $B_{ik}$,
and let $\Lambda$ be the sum of uniform twists with these weights.
It is a sum of exactly $L$ roots of unity.  The first twist moment of
$I(X)$ vanishes for nonunit $X$ \cite[Theorem 2.5]{Ostrik3}, so, putting
$w=5U$, the five equations combine as
\[
 w\Lambda j+t\sum_{i,k}\Theta_{ik}B_{ik}
 =-mzj-t(4mj+\delta).
\]
Pairing with $\delta$ gives $\sum\Theta_{ik}=-2$; pairing with $j$ gives
\[
 w\Lambda=-mz-2mt,\qquad
 \Lambda=-\frac{m^2+m\sqrt n}{2}.
\]

For completeness the required root-of-unity bound holds for all orders.
For $K=\QQ(\sqrt{q_0})$ with positive fundamental discriminant $f$, let
$c=\sqrt{f/q_0}$.  The coefficient of $\sqrt{q_0}$ in the normalized trace
to $K$ of a single root of unity has absolute value at most
$c/\varphi(f)$.  If the root's cyclotomic field does not contain $K$,
the trace is rational.  Otherwise its order is divisible by $f$.
Averaging over the kernel of the primitive quadratic character expresses
the irrational part as an induced quadratic Gauss sum divided by
$\varphi(N)$.  This sum is zero unless the extra factor $N/f$ is squarefree
and coprime to $f$; in the nonzero case its magnitude is $\sqrt f$.
The formula follows directly by Chinese remaindering: excess exponents
at primes of $f$ cancel over lifts, and the coprime factor contributes its
Ramanujan sum.  As $\varphi(N)\ge\varphi(f)$, the stated bound follows.
Normalized trace and the triangle inequality then require
\[
 L\ge\frac{\varphi(f)}c\left|[\sqrt{q_0}]\Lambda\right|
   =\frac{\varphi(f)mb}{2c}.
\]
For $m=7$, $n=q_0=f=1229$, $b=c=1$, and $\varphi(f)=1228$.
The forced sum of 590 roots is $-(49+7\sqrt{1229})/2$, but its required
length is at least $1228\cdot7/2=4298>590$.  This contradiction proves the
pseudo-unitary exclusion.  The separate complete signed trace-ideal
certificate supplies a pseudo-unitary Galois conjugate of any hypothetical
categorification, proving the general exclusion.
\end{proof}

The executable certificate is tensor-based and invariant under relabelling
of the five nonunit objects.  Its normal and optimized-Python replay
reconstructs every low and high column, rejects missing branches and
altered arithmetic.  The frozen
40-input continuation scan applies the criterion to every input; only
$\RR{7}{3582}$ and $\RR{8}{6520}$ match this tensor family, and both are
excluded in the pseudo-unitary scope. The independent trace-ideal
reduction makes the latter exclusion unrestricted. The scan
of all $16612$ tensors through multiplicity eight checks the parameter-three through parameter-seven
members uniformly:
\[
\begin{array}{c|r|r}
\text{tensor}&\text{forced root count}&\text{Gauss lower bound}\\\hline
\RR{4}{803}&110&342\\
\RR{5}{1500}&194&400\\
\RR{6}{2554}&302&1440\\
\RR{7}{3582}&434&1344\\
\RR{8}{6520}&590&4298
\end{array}
\]
The table records a separate application of the same necessary criterion
to all five rings.
